\chapter{Jev and ``System One'' models}
\label{ch:jev}

\section{What Jev does, in one line}
\label{sec:jev-oneline}

Jev reads a text or JSON \term{state} and answers typed, closed questions about it,
returning probabilities and (for two of the three question types) a confidence
statistic. It does not generate text.\DOC{}~\citep{typesafe_systemone}

TypeSafe calls it a \term{System One} model, after Kahneman's fast, intuitive ``System~1'':
it makes quick bounded judgments, not deliberate multi-step reasoning.

\section{The three primitives}

\begin{figure}[h]\centering
\begin{tikzpicture}[node distance=5mm and 10mm, every node/.style={font=\scriptsize}]
  \node[model] (choice) {\textbf{Choice}\\pick 1 of $\le$255\\options};
  \node[model, right=of choice] (score) {\textbf{Score}\\position on a\\2--10 level scale};
  \node[model, right=of score] (noul) {\textbf{Noul}\\probability a\\statement is true};
  \node[infra, below=6mm of choice] (c2) {returns choice,\\probabilities,\\confidence};
  \node[infra, below=6mm of score] (s2) {returns score,\\legend, probs,\\confidence};
  \node[infra, below=6mm of noul] (n2) {returns a single\\probability\\(no confidence)};
  \draw[flow] (choice)--(c2); \draw[flow] (score)--(s2); \draw[flow] (noul)--(n2);
\end{tikzpicture}
\caption{The three Jev primitives and what each returns.\datalabel{Reported by TypeSafe docs}~\citep{typesafe_primitives}}
\label{fig:primitives}
\end{figure}

\begin{description}[leftmargin=1.4cm,style=nextline]
  \item[Choice] one of up to 255 unordered options; returns the pick, a probability per
    option, and a confidence.\DOC{}
  \item[Score] a position on an ordered 2--10 level scale; the score is a
    probability-weighted mean of level numbers. Numerical interpolation is weakly
    calibrated, so use scores for thresholds, not arithmetic.\DOC{}
  \item[Noul] the probability that a statement is true; no separate confidence field.\DOC{}
\end{description}

\section{Anatomy of a request and response}

You send one \emph{state} and a set of independent \emph{questions}; they are evaluated in
parallel and cannot see each other's answers.\DOC{} Dependent questions need a second
request.

\begin{lstlisting}[language=jsonc,caption={A Choice request and its response (schematic).},label={lst:jevreq}]
// request
{ "state": { "task": "book a flight", "last_actions": ["search_flights"] },
  "model": "jev-1.13.0",
  "questions": [ { "id": "next_tool", "type": "choice",
                   "options": ["select_seat","pay","search_flights","none_of_the_above"] } ] }
// response
{ "answers": { "next_tool": { "choice": "select_seat",
      "probabilities": {"select_seat":0.71,"pay":0.20,"search_flights":0.08,
                        "none_of_the_above":0.01},
      "confidence": 0.68 } },
  "model": "jev-1.13.0", "usage": {"input_tokens": 41} }
\end{lstlisting}

\begin{worked}
Cost is dominated by input tokens: at the list price of \$0.042 per million input tokens
with free output, 20{,}000 decisions of about 1{,}500 input tokens each cost roughly
\$1.26.\DOC{}~\citep{typesafe_models} The limits that shape the SDK: at most 255 Choice
options, 2--10 Score levels, and 64k tokens per request (32k for state plus the longest
question).\DOC{}
\end{worked}

\section{The nine documented failure modes}

TypeSafe publishes a ``jaggedness'' page listing where Jev fails and how to work around
it.\DOC{}~\citep{typesafe_jaggedness} These are load-bearing for the design.

\begin{figure}[h]\centering
\small
\begin{tabular}{@{}p{4.6cm}p{7.4cm}@{}}
\toprule
Failure mode & Consequence for the design \\
\midrule
1. Literal reading & Question wording becomes a versioned artifact. \\
2. Math / counting & Keep arithmetic and counts in code. \\
3. Date/time comparison & Extract components; compute dates in code. \\
4. Indirection / multi-hop & Excludes planning-level decisions. \\
5. Large irrelevant state & Send projections, never full history. \\
6. Adversarial content & Cannot be the security boundary. \\
7. Contradictory criteria & Criteria design discipline. \\
8. No cross-primitive invariants & Do not transfer thresholds between primitives. \\
9. No generation & The LLM writes all text and arguments. \\
\bottomrule
\end{tabular}
\caption{The nine documented failure modes and what each forces in the architecture.\datalabel{Reported by TypeSafe docs}~\citep{typesafe_jaggedness}}
\label{fig:failuremodes}
\end{figure}

\begin{whyhere}
Every one of these nine drives a concrete rule later: projections (mode~5), deterministic
verifiers (modes~2--3), advisory-only authority (mode~6), frozen question wording
(mode~1), and per-primitive thresholds (mode~8). The design is essentially a disciplined
response to this table.
\end{whyhere}

\begin{checkbox}
Jev returns a probability of \(0.99\) that a shell command is safe. Give two reasons from
this chapter why code, not Jev, must still decide whether to run it. (Appendix~\ref{app:answers}.)
\end{checkbox}
